\begin{abstract}
本道题目数据集为 50,000 条有关电动汽车的消费者记录，三个子任务分别为对电动汽车采用意愿、家庭充电安装条件和里程焦虑敏感人群的估计和识别。

本题难点在于数据庞大，除开数据条数较多外，二十多列数据维度也为数据分析任务提出了挑战。
在数据预处理阶段，本文识别出部分维度存在缺失，在补足后又进行了缩尾和对数变换以提升数据质量。最终得到了一个 34 乘以 50,000 的数据集。
在模型选择上，考虑到数据集中有较多的分类变量，而分析任务均为归类识别，本文采取了 CatBoost 模型而非 LightBGM 或 XGBoost 模型。

在子任务一中，本文额外添加了 8 个交互特征，并使用 CatBoost 模型在测试集中得到了 86.33\% 的准确率。
在子任务二中，本文尝试了包括原始特征（赛题推荐特征）、诊断上界、跨任务堆叠三层级对照设计，甚至是随机森林算法均无法突破 0.65 左右的准确率。通过 EDV，本文认为，预测效果较差主要是由于原始数据中特征变量区分度较低。
在子任务三中，由于题目要求以 5 为边界将答案二值化，而大量分布在 5 附近的数据特征分布高度重叠，要在这个区间做好预测，需要通过 EDV 构造出精细的交互特征变量，以及本文特别引入了主方案-消融方案对照模型，并尤其注重位于 5 附近的边界样本分析，以探究模型对集中分布于中间的边界样本的预测水平。
最终，本文给出并实现了一套统一的 CatBoost 建模框架、跨任务无泄漏的堆叠设计。而任务二的失败则是一个从 EDV 到负结果报告的反面教材。

本文所有代码和中间文件均在 https://github.com/blackwhitetony/innoCupTaskB 仓库中开源。

\textbf{关键词：}CatBoost 电动汽车 分类预测 EDV
\end{abstract}
